\(X\ge0\) 且 \(\E[X]=5\)。Markov 不等式给出 \(P(X\ge 20)\le 5/20=0.25\)。若还知道 \(\Var(X)=4\)，Chebyshev 给出 \(P(|X-5|\ge 10)\le 4/100=0.04\)。界很松不表示事件危险，只表示仅凭一阶矩或二阶矩无法排除它——这就是不等式与精确分布计算的根本区别。

完整分布往往很难求。许多独立随机量相加时，逐次卷积越来越繁琐；面对尾部事件时，我们可能只知道均值、方差或取值范围。概率论因此发展出两类压缩工具：

\begin{itemize}
\tightlist
\item
  变换把分布编码成函数，并把独立和的卷积变成乘法；
\item
  不等式用有限信息给出尾部或坏事件的上界。
\end{itemize}

本单元依照``可用信息逐渐增多，界逐渐变强''的顺序阅读：

\begin{enumerate}
\tightlist
\item
  生成函数、矩母函数与特征函数：比较三种变换的适用对象、存在条件和求和规则。
\item
  从 Union Bound 到 Markov 与 Chebyshev：只用事件分解、均值或方差控制概率。
\item
  Jensen 不等式与随机凸性：平均与非线性运算为何不能随意交换。
\item
  Chernoff 方法与 Hoeffding 集中界：利用指数矩和独立性得到指数衰减的尾界。
\end{enumerate}

概率界给的是上界，不是精确概率。界很松并不意味着事件常发生，只说明现有信息不足以排除它；界很强也必须建立在相应独立性、矩存在或有界性条件上。

\subsection{一个具体算例}

抛 100 枚公平硬币，正面数 \(X\sim\Binomial(100,0.5)\)。Chebyshev：\(P(|X-50|\ge 10)\le 25/100=0.25\)。Chernoff/Hoeffding：\(P(|X-50|\ge 10)\le 2e^{-2\cdot 100/100}=2e^{-2}\approx 0.27\)（Hoeffding 对一般有界情形），或 \(2e^{-2\cdot 10^2/100}=2e^{-2}\approx 0.27\)——好不到哪去因为方差太大。改用 \(n=10000\)：Chebyshev 给出 \(P(|X-5000|\ge 100)\le 2500/10000=0.25\)，Hoeffding 给出 \(2e^{-2\cdot 100^2/10000}=2e^{-2}\approx 0.27\)——两者在 \(n\) 不大时都不精确；\(n\) 极大时 Hoeffding 才给出指数衰减。

Jensen 不等式实例：\(\E[X^2]\ge(\E[X])^2\)（取 \(\varphi(x)=x^2\)，凸函数），等号成立当且仅当 \(X\) 几乎处处为常数。因此方差 \(\Var(X)=\E[X^2]-(\E[X])^2\ge 0\) 是 Jensen 的直接推论。

\subsection{怎么读这一章}

若目标是识别或组合分布，先读第一篇，并在使用 PGF、MGF、特征函数前分别检查取值类型、存在域与独立性。若目标只是控制坏事件，按可用信息选择：只有事件概率用 Union Bound，只有非负均值用 Markov，有有限方差用 Chebyshev 或 Cantelli，有独立有界结构再用 Hoeffding；Jensen 则负责判断非线性运算与平均的先后代价。

\subsection{本章篇目}

以下篇目按概念依赖排列；若从具体问题进入，也可以先打开最接近当前困惑的一篇，再沿篇内链接回补。

\begin{itemize}
\tightlist
\item \hyperref[sec:05-Probability/07-Transforms-and-Bounds/01]{``生成函数、矩母函数与特征函数''}；
\item \hyperref[sec:05-Probability/07-Transforms-and-Bounds/02]{``从 Union Bound 到 Markov 与 Chebyshev''}；
\item \hyperref[sec:05-Probability/07-Transforms-and-Bounds/03]{``Jensen 不等式与随机凸性''}；
\item \hyperref[sec:05-Probability/07-Transforms-and-Bounds/04]{``Chernoff 方法与 Hoeffding 集中界''}。
\end{itemize}
